% ── SECTION: COHERENCE-BUDGET ALLOCATION (~v1.0 new content) ────────────────
%   v1.0 contribution: the allocation theorem, scheduling algorithm,
%   and multi-layer fidelity benchmark.

\section{Coherence-Budget Allocation for Tunable Exchange Schedules}
\label{sec:alloc}

The per-gate fidelity comparison of Sec.~\ref{sec:sim} establishes
that, on hardware with directly calibratable Heisenberg exchange,
$\URA(\theta)$ at $\theta < \pi/4$ accumulates less coherence-limited
error than a full maximally entangling primitive at equal entangling
capability requirement.
That observation, however, says nothing about how to \emph{allocate}
entanglement across the gate layers of a circuit whose total
entanglement target is fixed.
This section turns the per-gate fidelity statement into a layerwise
optimization: given a coherence budget and an entanglement target,
how should a noise-aware compiler distribute exchange angles across
the circuit?

\subsection{Relational Coherence Cost}
\label{subsec:cost}

I define the per-layer cost of a tunable exchange gate
$\URA(\theta_\ell)$ as the local pointer-basis coherence consumed by
the gate:
%
\begin{equation}
\label{eq:layer-cost}
C_{\RCF}^{(\ell)} \;\equiv\; \bar\alpha_\ell^{\rm pre} - \bar\alpha_\ell^{\rm post}
\;=\; \bar\alpha_\ell^{\rm pre}\bigl(1 - e^{-\gamma_\ell \kappa\, \theta_\ell}\bigr),
\end{equation}
%
where $\gamma_\ell$ is the local pointer-basis dephasing rate active
during layer $\ell$ (which may vary across layers due to hardware
heterogeneity), $t_\ell = \kappa\, \theta_\ell$ is the exchange-pulse
duration on the substrate~\cite{Krantz2019}, and the exponential
form follows from
Proposition~\ref{prop:determination} in its uniform-rate Markovian
regime.
The total circuit cost is
%
\begin{equation}
\label{eq:circuit-cost}
C_{\RCF}(\mathcal{U}) \;=\; \sum_{\ell=1}^{L} C_{\RCF}^{(\ell)}.
\end{equation}
%
To leading order in $\gamma_\ell \kappa\, \theta_\ell \ll 1$
(the small-cost regime in which Eq.~\eqref{eq:peff}'s
gate-time-proportional model itself is justified, Sec.~\ref{sec:sim}),
%
\begin{equation}
\label{eq:linear-cost}
C_{\RCF}(\mathcal{U}) \;\approx\;
  \sum_{\ell=1}^{L} \tilde\gamma_\ell\, \theta_\ell,
  \qquad \tilde\gamma_\ell \equiv \bar\alpha_\ell^{\rm pre}\, \gamma_\ell\, \kappa.
\end{equation}
%
The effective rate $\tilde\gamma_\ell$ folds the substrate's exchange
time-constant and the layer's residual pre-gate coherence into a
single per-layer cost coefficient.

\subsection{Proposition~V: Exchange-Angle Coherence Cost}
\label{subsec:prop5}

The cost per layer admits a sharp closed form for the
directly-calibrated exchange-gate case treated throughout this paper.

\begin{proposition}[Exchange-angle coherence cost, Markovian regime]
\label{prop:exchange-cost}
Under uniform Markovian pointer-basis dephasing at rate $\gamma$ and
direct exchange-gate duration $t(\theta) = t_{\pi/4}\,\theta/(\pi/4)$
(Sec.~\ref{sec:sim}), the layerwise scalar-affinity loss induced by
$\URA(\theta)$ on a layer with input affinity $\bar\alpha_{\rm in}$ is
%
\begin{equation}
\label{eq:exchange-cost}
\Delta\bar\alpha(\theta)
  \;=\; \bar\alpha_{\rm in}\!\left[\,
        1 - \exp\!\left(-\gamma\,t_{\pi/4}\,\frac{\theta}{\pi/4}\right)\right].
\end{equation}
%
In the small-error regime $\gamma\, t_{\pi/4}\,\theta/(\pi/4) \ll 1$,
%
\begin{equation}
\label{eq:exchange-cost-linear}
\Delta\bar\alpha(\theta)
  \;\approx\; \gamma\,t_{\pi/4}\,\frac{\theta}{\pi/4}\;\bar\alpha_{\rm in}.
\end{equation}
%
The per-layer coherence cost is therefore first-order linear in the
exchange angle $\theta$.
\end{proposition}

\begin{proof}
By Proposition~\ref{prop:determination} (Markovian regime), the scalar
affinity decays as
$\bar\alpha(t) = \bar\alpha_{\rm in}\,e^{-\gamma t}$.
Substituting $t = t_{\pi/4}\,\theta/(\pi/4)$ gives
$\bar\alpha^{\rm post} = \bar\alpha_{\rm in}\,\exp(-\gamma\,t_{\pi/4}\,\theta/(\pi/4))$,
and $\Delta\bar\alpha \equiv \bar\alpha_{\rm in} - \bar\alpha^{\rm post}$
gives Eq.~\eqref{eq:exchange-cost}.
The linearization Eq.~\eqref{eq:exchange-cost-linear} follows from
$1 - e^{-x} = x + O(x^2)$. $\square$
\end{proof}

\begin{remark}[Linear cost matches Eq.~\eqref{eq:linear-cost}]
The leading-order form Eq.~\eqref{eq:exchange-cost-linear} reproduces
the per-layer term in Eq.~\eqref{eq:linear-cost} with effective rate
$\tilde\gamma_\ell = \bar\alpha_\ell^{\rm pre}\,\gamma_\ell\,t_{\pi/4}/(\pi/4)$;
the constant $\kappa = t_{\pi/4}/(\pi/4)$ is the substrate's
exchange-time scale per unit angle.
The allocation problem of Sec.~\ref{subsec:problem} below is therefore
the layer-aggregated form of
Proposition~\ref{prop:exchange-cost}.
\end{remark}

\subsection{Corollary: Entanglement Efficiency}
\label{subsec:eff}

Combining Proposition~\ref{prop:exchange-cost} with the
entanglement-capability formula
$\mathcal{E}(\theta) = |\sin 2\theta|$ (Sec.~\ref{subsec:entcap})
yields the per-gate entanglement-per-coherence-cost ratio.

\begin{corollary}[Entanglement efficiency]
\label{cor:eta}
In the small-error regime of
Proposition~\ref{prop:exchange-cost}, the entanglement
generated per unit of pointer-basis coherence spent on a single
exchange gate is
%
\begin{equation}
\label{eq:eta}
\eta(\theta) \;\equiv\;
  \frac{\mathcal{E}(\theta)}{\Delta\bar\alpha(\theta)}
  \;\propto\; \frac{\sin 2\theta}{\theta},
\end{equation}
%
suppressing the constant factor
$\bigl(\gamma\,t_{\pi/4}\,\bar\alpha_{\rm in}/(\pi/4)\bigr)^{-1}$.
The ratio $\sin(2\theta)/\theta$ is strictly decreasing on
$(0, \pi/4]$, with $\lim_{\theta\to 0^+}\sin(2\theta)/\theta = 2$
and $\sin(\pi/2)/(\pi/4) = 4/\pi \approx 1.273$.
\end{corollary}

\begin{remark}[Efficiency vs.\ capacity]
Corollary~\ref{cor:eta} formalizes the tension exploited by the
allocation theorem.
Small-$\theta$ gates are most \emph{efficient} per unit coherence
spent (high $\eta$), but produce small absolute entanglement
$\mathcal{E}(\theta) \approx 2\theta$ per gate.
Large-$\theta$ gates are less efficient ($\eta \to 4/\pi$ at the
$\sqrt{\textsc{swap}}$ point) but achieve more capacity per gate.
The Coherence-Budget Allocation Theorem
(Sec.~\ref{subsec:thm-alloc}) resolves this tension by distributing
entanglement across layers when $L > 1$ allows, and concentrates it
when noise heterogeneity forces a tradeoff.
\end{remark}

\subsection{The Allocation Problem}
\label{subsec:problem}

For a circuit with required total entangling capability
$E_{\rm target}$ accumulated across $L$ exchange layers, each
contributing $\mathcal{E}(\theta_\ell) = \sin(2\theta_\ell)$
(Sec.~\ref{subsec:entcap}; I restrict to the entangling half
$\theta_\ell \in [0, \pi/4]$ on which $\sin(2\theta_\ell)$ is
increasing), the coherence-optimal schedule solves
%
\begin{equation}
\label{eq:alloc-problem}
\begin{aligned}
\text{minimize}\quad & C_{\RCF}(\boldsymbol\theta) \;=\;
   \sum_{\ell=1}^L \tilde\gamma_\ell\, \theta_\ell \\
\text{subject to}\quad & \sum_{\ell=1}^L \sin(2\theta_\ell) \;\geq\; E_{\rm target}, \\
& 0 \leq \theta_\ell \leq \pi/4, \qquad \ell = 1,\dots, L.
\end{aligned}
\end{equation}
%
The objective is linear (hence convex) and the constraint set is
convex: $\sum_\ell \sin(2\theta_\ell)$ is concave on
$[0,\pi/4]^L$ (second derivative $-4\sin 2\theta_\ell \leq 0$
component-wise), so the superlevel set
$\{\boldsymbol\theta : \sum_\ell \sin(2\theta_\ell) \geq E_{\rm target}\}$
is convex.
Minimizing a linear functional on a convex set yields a unique
minimizer (interior) or a minimizing face (boundary).

\subsection{The Coherence-Budget Allocation Theorem}
\label{subsec:thm-alloc}

\begin{theorem}[Coherence-Budget Allocation]
\label{thm:alloc}
Let $\tilde\gamma_1,\dots,\tilde\gamma_L > 0$ be the effective
per-layer dephasing rates of Eq.~\eqref{eq:linear-cost} and let
$E_{\rm target} \in (0, L)$.
The unique solution of the convex program
Eq.~\eqref{eq:alloc-problem} is given layerwise by
%
\begin{equation}
\label{eq:alloc-solution}
\theta_\ell^{\!*} \;=\;
  \frac{1}{2}\,\arccos\!\Bigl(
    \min\!\Bigl(1,\;\tfrac{\tilde\gamma_\ell}{2\lambda^{\!*}}\Bigr)
  \Bigr),
  \qquad \ell = 1,\dots,L,
\end{equation}
%
where the shadow price $\lambda^{\!*} > 0$ is the unique
solution of the budget constraint
%
\begin{equation}
\label{eq:lambda-eq}
\sum_{\ell=1}^L \sin\!\bigl(2\theta_\ell^{\!*}(\lambda^{\!*})\bigr)
\;=\; E_{\rm target}.
\end{equation}
\end{theorem}

\begin{proof}
The Lagrangian is
$\mathcal{L}(\boldsymbol\theta,\lambda,\boldsymbol\mu,\boldsymbol\nu)
= \sum_\ell \tilde\gamma_\ell\,\theta_\ell
- \lambda\!\bigl(\sum_\ell \sin 2\theta_\ell - E_{\rm target}\bigr)
- \sum_\ell \mu_\ell\,\theta_\ell
+ \sum_\ell \nu_\ell\,(\theta_\ell - \tfrac{\pi}{4})$
with $\lambda,\mu_\ell,\nu_\ell \geq 0$.
Stationarity $\partial\mathcal{L}/\partial\theta_\ell = 0$ gives
%
\begin{equation*}
\tilde\gamma_\ell - 2\lambda \cos 2\theta_\ell - \mu_\ell + \nu_\ell = 0.
\end{equation*}
%
For an interior layer ($0 < \theta_\ell^* < \pi/4$, complementary
slackness $\mu_\ell = \nu_\ell = 0$) this reduces to
$\cos 2\theta_\ell^* = \tilde\gamma_\ell/(2\lambda)$,
giving Eq.~\eqref{eq:alloc-solution} when the argument is in $[0,1]$.
For the lower boundary $\theta_\ell^* = 0$ ($\nu_\ell = 0$,
$\mu_\ell \geq 0$), stationarity requires
$\tilde\gamma_\ell - 2\lambda \geq 0$, i.e., $\tilde\gamma_\ell \geq 2\lambda$,
which is exactly the condition $\tilde\gamma_\ell/(2\lambda) \geq 1$
in Eq.~\eqref{eq:alloc-solution} (the $\min(\cdot)$ clip).
The upper boundary $\theta_\ell^* = \pi/4$ requires $\tilde\gamma_\ell = 0$,
excluded by hypothesis.

Existence and uniqueness of $\lambda^{\!*}$:
$\sin 2\theta_\ell^*(\lambda)$ is monotone non-decreasing and
continuous in $\lambda$ on $(0,\infty)$, with limits $0$ as $\lambda\to 0^+$
and $1$ as $\lambda\to\infty$.
Hence $\Sigma(\lambda) \equiv \sum_\ell \sin 2\theta_\ell^*(\lambda)$
is continuous and monotone, mapping $(0,\infty)$ onto $(0, L)$;
$\Sigma(\lambda^{\!*}) = E_{\rm target} \in (0,L)$ has a unique
solution.
Global optimality follows from the convexity of the program. $\square$
\end{proof}

\subsection{Three Falsifiable Predictions}
\label{subsec:predictions}

Three predictions follow immediately from Eq.~\eqref{eq:alloc-solution}
and are testable in numerical simulation.

\textbf{P1 (Distribution beats concentration, uniform-$\tilde\gamma$).}
For $\tilde\gamma_\ell = \tilde\gamma$ uniform, the optimal schedule is
$\theta_\ell^* = \tfrac{1}{2}\arcsin(E_{\rm target}/L)$ on all layers,
with total cost
$C^* = \tilde\gamma\,L\cdot \tfrac{1}{2}\arcsin(E_{\rm target}/L)$.
By the strict concavity of $\arcsin$,
$L \arcsin(E_{\rm target}/L) < \arcsin(E_{\rm target})$ for $L > 1$
and $E_{\rm target} \in (0,1]$;
hence the optimal distributed schedule strictly beats the cost of
implementing $E_{\rm target}$ in a single partial entangler.

\textbf{P2 (Quieter qubits receive more entanglement).}
For heterogeneous $\tilde\gamma_\ell$, layers with smaller
$\tilde\gamma_\ell$ have smaller $\tilde\gamma_\ell/(2\lambda^{\!*})$,
hence smaller $\cos 2\theta_\ell^*$, hence larger $\theta_\ell^*$.
The optimal compiler routes entangling angle toward quieter substrates.

\textbf{P3 (Skip threshold for noisy qubits).}
Layers with $\tilde\gamma_\ell \geq 2\lambda^{\!*}$ are clipped to
$\theta_\ell^* = 0$: the optimum spends no coherence on layers whose
local cost exceeds the shadow price of the global budget constraint,
provided quieter alternatives suffice to meet $E_{\rm target}$.

P1 is a no-heterogeneity baseline statement;
P2 and P3 are the practically distinguishing predictions on
realistic hardware whose dephasing rates vary across qubit pairs.

\subsection{Scheduling Algorithm}
\label{subsec:algorithm}

The shadow-price equation Eq.~\eqref{eq:lambda-eq} is a monotone
one-dimensional root-finding problem and admits a Brent-method
solution in milliseconds for $L \lesssim 10^3$.
The full algorithm is:

\begin{enumerate}
  \item Input: layer rates $\{\tilde\gamma_\ell\}_{\ell=1}^L$,
        entanglement target $E_{\rm target}$.
  \item Bracket $\lambda$: find $[\lambda_{\rm lo}, \lambda_{\rm hi}]$
        such that $\Sigma(\lambda_{\rm lo}) < E_{\rm target} < \Sigma(\lambda_{\rm hi})$.
  \item Solve $\Sigma(\lambda^{\!*}) = E_{\rm target}$ by Brent's method.
  \item Return $\theta_\ell^* = \tfrac{1}{2}\arccos(\min(1,\tilde\gamma_\ell/(2\lambda^{\!*})))$.
\end{enumerate}

The reference implementation (\texttt{figures/rcb\_optimizer.py})
runs in $O(L \log(1/\epsilon))$ time for tolerance $\epsilon$ and is
included with this paper.

\subsection{Multi-Layer Benchmark}
\label{subsec:benchmark-alloc}

I test Theorem~\ref{thm:alloc} against two naive scheduling
baselines on a multi-layer 2-qubit circuit with per-layer independent
dephasing.

\textbf{Setup.}
A circuit of $L = 8$ layers applies $\URA(\theta_\ell)$ followed by
independent per-qubit pointer-basis dephasing at rate $\gamma_\ell$
for duration $t_\ell = \kappa\,\theta_\ell$, starting from
$\rho_0 = \ket{+0}\!\bra{+0}$.
For each entanglement target $E_{\rm target} \in [0.25, 3.5]$, I
compare three strategies:
\begin{itemize}
  \item \textbf{RCB-optimal} schedule per Theorem~\ref{thm:alloc};
  \item \textbf{Uniform $\theta$}: $\theta_\ell = \tfrac{1}{2}\arcsin(E_{\rm target}/L)$
        on all layers (the optimal schedule when $\tilde\gamma_\ell$ is uniform);
  \item \textbf{Concentrated}: place full entanglers
        $\theta_\ell = \pi/4$ on the $\lfloor E_{\rm target}\rfloor$
        quietest layers, residual on the next-quietest.
\end{itemize}
For each strategy I compute the final reduced-state fidelity
$F = \mathrm{Tr}[\sqrt{\sqrt{\rho_{\rm noisy}}\,\rho_{\rm ideal}\,\sqrt{\rho_{\rm noisy}}}]^2$
against the noiseless target unitary.

\textbf{Results.}
Figure~\ref{fig:rcb-sweep} shows fidelity (top row) and coherence cost
(bottom row) as $E_{\rm target}$ varies, for three $\gamma$ regimes:
uniform low ($\gamma = 0.05$), uniform high ($\gamma = 0.15$), and
heterogeneous ($\gamma \in [0.02, 0.30]$).
The RCB-optimal schedule meets or exceeds both baselines for every
$E_{\rm target}$ tested, with the largest fidelity advantage in the
high-$\gamma$ regime (e.g., $\Delta F = +0.040$ over concentrated at
$E_{\rm target} = 1$, $\gamma = 0.15$) and a substantial advantage
over uniform-$\theta$ in the heterogeneous regime
($\Delta F = +0.025$ at $E_{\rm target} = 1$).
At representative targets:

\begin{center}
\small
\begin{tabular}{lccc}
\toprule
Regime & $\Delta F_{\rm opt - conc}$ at $E\!=\!1$ & at $E\!=\!1.5$ & at $E\!=\!2$ \\
\midrule
Uniform low ($\gamma\!=\!0.05$) & $+0.014$ & $+0.014$ & $+0.013$ \\
Uniform high ($\gamma\!=\!0.15$) & $\mathbf{+0.040}$ & $+0.037$ & $+0.033$ \\
Heterogeneous & $+0.005$ & $+0.004$ & $+0.001$ \\
\bottomrule
\end{tabular}
\end{center}

\begin{figure*}[htbp]
\centering
\includegraphics[width=0.9\textwidth]{fig3_rcb_sweep}
\caption{
  Coherence-Budget Allocation benchmark on an $L\!=\!8$ layer
  two-qubit circuit with per-layer independent pointer-basis dephasing.
  \textbf{Top row:} final-state fidelity vs noiseless target, as
  a function of entanglement target $E_{\rm target}$, for the RCB-optimal
  schedule (Eq.~\eqref{eq:alloc-solution}, solid),
  uniform-$\theta$ baseline (dashed), and concentrated full-entangler
  baseline (dotted).
  \textbf{Bottom row:} coherence cost
  $\sum_\ell \tilde\gamma_\ell\,\theta_\ell$ for the same schedules.
  \textbf{Columns:} uniform low noise
  ($\gamma\!=\!0.05$), uniform high noise ($\gamma\!=\!0.15$), and
  heterogeneous noise ($\gamma \in [0.02, 0.30]$).
  The RCB-optimal schedule equals or exceeds both baselines in
  fidelity across all $E_{\rm target}$; the largest advantage appears
  in the high-$\gamma$ regime (Prediction~P1) and in the
  heterogeneous regime against uniform $\theta$ (Prediction~P2).
  Generated by \texttt{figures/benchmark\_sweep.py}.
}
\label{fig:rcb-sweep}
\end{figure*}

P1 is visible in the uniform-$\gamma$ columns: distribution always
beats concentration.
P2 is visible in the heterogeneous column: RCB-optimal beats
uniform-$\theta$ by routing entanglement onto the
$\gamma\!=\!0.02$ layers and zero onto the $\gamma\!=\!0.30$ layer
(P3 in action).
The benchmark validates Theorem~\ref{thm:alloc} as a per-circuit,
not merely per-gate, statement of the coherence-budget framing.

\subsection{Scope of Theorem~\ref{thm:alloc}}
\label{subsec:alloc-scope}

The allocation theorem holds under the assumptions made explicit in
Sec.~\ref{sec:sim}: (i)~Markovian pure-dephasing dynamics in the
einselected pointer basis, with rates that are constant within each
gate window;
(ii)~gate-time-proportional coherence-limited error (small-$p$
leading-order linearization of Eq.~\eqref{eq:peff});
(iii)~direct exchange calibration so that $\URA(\theta)$ is realized
as a single tunable pulse rather than decomposed into fixed native
gates.
Where (iii) fails, the constant $\kappa$ relating $\theta$ to
duration in Eq.~\eqref{eq:layer-cost} is replaced by a
decomposition-depth--dependent factor; the optimization problem and
its solution form are unchanged but the numerical advantage shrinks
as in Sec.~\ref{sec:disc}.
Coherent calibration errors, leakage to non-computational levels, and
crosstalk are not modelled by Eq.~\eqref{eq:peff} and can shift the
optimum; their incorporation as additional per-layer cost terms is
straightforward and preserves the convex-program structure.
